\documentclass{situationist-preprint}

\title{The Role of Notation}
\msid{MTH-03-b}
\msfamily{MTH-03}
\msdiscipline{Philosophy of Mathematics}
\mstype{Companion manuscript}
\mskeywords{notation; invention; discovery; positional numerals; Leibniz; mathematical practice}

\begin{document}
\maketitle

\begin{abstract}
MTH-03-a argues that mathematical truths are discovered while notations are invented. Notation nonetheless seems to do more than an arbitrary choice should: good notation makes results findable that bad notation hides, and some notations seem to carry truths within them. We argue that notation has an epistemic role and no constitutive one. It determines which truths are accessible, and at what cost, but not which are true; the persistence of results across notations, which MTH-03-a treats as a mark of discovery, is the evidence.
\end{abstract}

\section{Two roles}\label{s:roles}

A notation might matter to mathematics in one of two ways. It might be \emph{constitutive}, partly making true the truths expressed in it, or it might be \emph{epistemic}, affecting which truths we can find and recognize, and how easily. Proposition~3.1 of MTH-03-a classifies notations as inventions; the question is whether their evident importance threatens Theorem~3.3, that mathematical truths are discovered.

The importance is not in doubt. Positional decimal numerals turned multiplication from an art into a routine, and Leibniz's $dy/dx$ makes the chain rule look like the cancellation of fractions, so that it can be applied before it is understood \citep{Cajori1928}. Frege devised his concept-script because ordinary language hid the gaps in proofs \citep{Frege1879}.

\begin{proposition}\label{p:epi}
Notation has an epistemic role in mathematics and no constitutive role.
\end{proposition}

\begin{proof}
If a notation were constitutive of a truth, the truth would fail when the notation was replaced. By the invariance mark of Definition~2.2 of MTH-03-a, results survive changes of notation: the chain rule holds in Newton's fluxional notation and in the language of linear maps, and the infinitude of the primes holds in every numeral system. No result is known that holds in one notation and fails in an equivalent one. On the other hand, the cost of finding and checking results differs greatly across notations, as the examples show. So notation affects access and not truth.
\end{proof}

The epistemic role has an explanation. A notation is the perceptible material in which proofs are written and read; PHE-14 argues that the evidence of formal reasoning is founded on perceptual acts of reading and writing, while the truth of what is proved does not depend on perception. A good notation is one in which structural features of the subject matter are mirrored by visible features of the inscriptions, so that inspecting the inscription is a reliable way of tracking the structure. That is why the choice of notation, though free in the sense of Definition~2.1 of MTH-03-a, is not indifferent.

\begin{objection}[Suggestive notation]
Some notations lead to truths their inventors did not anticipate: Leibniz's differentials suggested the rules for integration by substitution before they were proved. If notation can lead us to truths, it is more than a tool.
\end{objection}

\begin{reply}
A notation leads to truths because it mirrors structure, and it misleads where it fails to: $dy/dx$ suggests that $(dy/dx)^2 = dy^2/dx^2$, which is not the second derivative. The suggestions of a notation must be checked against the subject matter, and are sometimes rejected. A notation whose suggestions were true by fiat would never mislead. That notation can mislead is evidence that it answers to something it does not make.
\end{reply}

\section{Conclusion}

Notation is invented and important. Its importance is epistemic: it controls access to mathematical truth, as an instrument controls access to a physical quantity, and it can be better or worse at the job. That is consistent with, and partly explained by, the thesis that mathematical truths are discovered.

\begin{thebibliography}{9}\small
\bibitem[Cajori(1928)]{Cajori1928} Cajori, F. (1928--1929). \emph{A History of Mathematical Notations}. 2 vols. Open Court.
\bibitem[Frege(1879)]{Frege1879} Frege, G. (1879). \emph{Begriffsschrift, eine der arithmetischen nachgebildete Formelsprache des reinen Denkens}. Louis Nebert.
\end{thebibliography}

\end{document}
